\section{Attention Residuals (AttnRes)}
\label{sec:annex-attnres}

This annex documents the integration of \emph{Attention Residuals}
(AttnRes) \citep{sun_attnres_2026}, a generalization of the residual
connection that replaces the fixed unit-coefficient residual sum with
\emph{learned softmax attention over depth}. The main body delegates
to this annex for the mathematical formulation, the four available
strategies (standard, no-residual, Full AttnRes, Block AttnRes), the
implementation wiring, the integration with mHC and Mixture-of-Depths,
and the reported results. The bibliography keys cited here resolve
against \texttt{docs/bibliography/}.

\subsection{Motivation: From a Fixed Sum to Depth-Wise Attention}

The standard residual connection \citep{he_deep_2016} propagates signals
across depth unmodified: with $x_l \in \mathbb{R}^{C}$ and a layer
function $F$,
\[
x_{l+1} = x_l + F(x_l, W_l),
\]
so that, recursively,
$x_L = x_l + \sum_{i=l}^{L-1} F(x_i, W_i)$. The identity term $x_l$
(the \emph{identity mapping}) lets information flow from shallow to deep
layers unchanged, which He et al.~\citep{he_deep_2016} identified as key
to stable training of very deep networks. Highway networks
\citep{srivastava_highway_2015} generalize this with learned element-wise
gates $\mathbf{g}_l \in [0, 1]^{C}$,
\[
x_{l+1} = (1 - \mathbf{g}_l) \odot x_l + \mathbf{g}_l \odot F(x_l, W_l),
\]
which interpolates between the identity path and the transformation
path with input-dependent weights.

\emph{Attention Residuals} \citep{sun_attnres_2026} take the next step:
instead of weighting a single compressed state $x_l$, layer $l$
\emph{attends} over all earlier layer outputs. Let
$v_0 = h_1$ (the token embedding) and $v_i = F(x_i, W_i)$ for $i \ge 1$.
AttnRes computes, for each layer $l$,
\[
q_l = w_l, \qquad
k_i = \mathrm{RMSNorm}(v_i), \qquad
\alpha_{i \to l} = \mathrm{softmax}_i(q_l^{\top} k_i), \qquad
x_{l+1} = \sum_{i=0}^{l} \alpha_{i \to l} \, v_i,
\]
where $w_l \in \mathbb{R}^{C}$ is a learned per-layer pseudo-query,
\emph{initialized to zero} per the paper. At zero-init the softmax is
uniform, so AttnRes degenerates to an equal-weight average of all
previous outputs --- matching the standard residual at step zero and
avoiding training volatility. The RMSNorm on the keys prevents layers
with naturally large magnitudes from dominating the softmax
(ablation in the paper: \S~5.3).

\subsection{Four Strategies}

The schema field \texttt{model.residuals.type} selects one of four
strategies, each implemented as a separate module under
\texttt{src/model/residuals/}.

\subsubsection{Standard Residual (\texttt{"standard"})}

The default backwards-compatible strategy:
\[
x_{l+1} = x_l + F(x_l, W_l).
\]
Stateless, no extra parameters, no extra compute. Implemented as a
no-op pass-through in \texttt{StandardResidual} that delegates the
merge to \texttt{HybridLayer}.

\subsubsection{No Residual (\texttt{"none"}, experimental)}

Drops the skip connection entirely:
\[
x_{l+1} = F(x_l, W_l).
\]
Stateless, no extra parameters. Useful only as an ablation probe; the
identity term is critical for deep stacks (paper \S~2.1) and training is
typically unstable without it. Implemented as \texttt{NoResidual}.

\subsubsection{Full Attention Residuals (\texttt{"full\_attn"}, paper \S~3.1)}

Each layer attends over \emph{all} previous layer outputs:
\[
x_{l+1} = \sum_{i=0}^{l} \alpha_{i \to l} \, v_i, \qquad
\alpha_{i \to l} = \mathrm{softmax}_i\!\big(w_l^{\top} \mathrm{RMSNorm}(v_i)\big).
\]
Adds $L \cdot C$ parameters ($L$ = logical depth = \texttt{num\_layers}
$\times$ \texttt{num\_loops}, $C$ = \texttt{hidden\_size}), one query
vector per layer. Memory and compute grow as $O(L \, d)$ per token; at
typical depth ($L < 1000$) this is negligible. Implemented as
\texttt{FullAttentionResidual}; supports
\texttt{attnres\_gradient\_checkpoint} to wrap the attention in
\texttt{torch.utils.checkpoint} for memory-constrained runs.

\subsubsection{Block Attention Residuals (\texttt{"block\_attn"}, paper \S~3.2)}

Partitions the $L$ logical layers into $N$ blocks of
$S = \lceil L / N \rceil$ layers each. Within a block the layer outputs
accumulate as a standard partial sum $b_n^i$; across blocks attention
is applied over the $N$ complete block representations and the running
partial sum. For layer $l$ in block $n$,
\[
x_{l+1} = \sum_{m=0}^{n} \alpha_{m \to l} \, v_m, \qquad
\alpha_{m \to l} = \mathrm{softmax}_m\!\big(w_l^{\top} \mathrm{RMSNorm}(v_m)\big),
\]
with $v_m = b_m$ for $m < n$ and $v_n = b_n^{i_l}$ (the current
intra-block partial sum). $N$ interpolates between Full AttnRes
($N = L$, one block per layer) and standard residuals ($N = 1$, the
embedding isolated as $b_0$). The paper's empirical sweet spot is
$N \approx 8$; this codebase defaults to $N = 8$ and validates that
$N \le L$. Implemented as \texttt{BlockAttentionResidual}, which also
follows Algorithm 1 of the paper (two-phase computation with online
softmax merging).

\subsection{Implementation in Frankenstein Transformer}

The \texttt{src/model/residuals/} package hosts one class per strategy
(\texttt{StandardResidual}, \texttt{NoResidual},
\texttt{FullAttentionResidual}, \texttt{BlockAttentionResidual}) plus
an abstract \texttt{ResidualBase} that defines the lifecycle hooks
\texttt{set\_streams}, \texttt{register\_state}, \texttt{reset\_state},
\texttt{forward}, and \texttt{finalize}. The
\texttt{build\_residual(config)} factory in
\texttt{src/model/residuals/factory.py} selects the right class from
\texttt{FrankensteinModelConfig}.

Wiring in \texttt{src/model/frankenstein\_encoder.py}:
\begin{itemize}[leftmargin=*]
    \item \texttt{FrankensteinEncoder} owns the
        \texttt{ResidualBase} module as \texttt{self.residual}; the
        factory is invoked once in \texttt{\_\_init\_\_}.
    \item Before the looped-depth loop the encoder calls
        \texttt{reset\_state} and, for AttnRes variants,
        \texttt{set\_embedding(x)} so the depth-wise attention knows
        the $v_0 = h_1$ source.
    \item After each \texttt{HybridLayer} call, if
        \texttt{self.residual.is\_attn\_res} is true, the encoder
        applies \texttt{x = self.residual(layer\_idx, x)} to overwrite
        the stream with the attended aggregation. For stateless
        strategies (\texttt{standard}, \texttt{none}) this step is
        skipped and \texttt{HybridLayer} continues to apply its own
        residual merge internally.
\end{itemize}

Wiring in \texttt{src/model/hybrid\_layer.py}:
\begin{itemize}[leftmargin=*]
    \item \texttt{HybridLayer} reads \texttt{residual\_type} from the
        config and dispatches the residual merge accordingly. For
        \texttt{standard} the merge is the original
        \texttt{x = residual + mixer(norm(x))}; for \texttt{none} the
        skip is dropped (\texttt{x = mixer(norm(x))}); for the AttnRes
        variants the layer performs the standard merge internally
        because the depth-wise attention is applied externally by the
        encoder.
    \item \texttt{\_forward\_dense\_mhc} (the mHC path) is unaffected
        because mHC has its own stream-expansion residual; combining
        AttnRes with mHC is handled by the \texttt{mhc\_stream\_mode}
        field of the residual module.
\end{itemize}

Schema (hierarchical, with the flat-key renames shown in parentheses)
and defaults:
\begin{itemize}[leftmargin=*]
    \item \texttt{model.residuals.type} (\texttt{residual\_type},
        enum \texttt{[standard, none, full\_attn, block\_attn]},
        default \texttt{standard}): selects the strategy.
    \item \texttt{model.residuals.full\_attn.init\_query\_zero}
        (\texttt{full\_attn\_init\_query\_zero}, bool, default
        \texttt{true}): zero-init the per-layer query vectors
        $\mathbf{w}_l$.
    \item \texttt{model.residuals.full\_attn.use\_rmsnorm\_keys}
        (\texttt{full\_attn\_use\_rmsnorm\_keys}, bool, default
        \texttt{true}): apply RMSNorm on keys before the dot product.
    \item \texttt{model.residuals.block\_attn.num\_blocks}
        (\texttt{block\_attn\_num\_blocks}, int $\ge 1$, default
        $8$): number of block representations $N$; validated
        $N \le L$.
    \item \texttt{model.residuals.block\_attn.init\_query\_zero}
        (\texttt{block\_attn\_init\_query\_zero}, bool, default
        \texttt{true}): zero-init the per-layer query vectors.
    \item \texttt{model.residuals.block\_attn.use\_rmsnorm\_keys}
        (\texttt{block\_attn\_use\_rmsnorm\_keys}, bool, default
        \texttt{true}): RMSNorm on keys.
    \item \texttt{model.residuals.mhc\_stream\_mode}
        (\texttt{attnres\_mhc\_stream\_mode}, enum
        \texttt{[independent, joint]}, default \texttt{independent}):
        how AttnRes interacts with mHC's n-stream residual.
    \item \texttt{model.residuals.gradient\_checkpoint}
        (\texttt{attnres\_gradient\_checkpoint}, bool, default
        \texttt{false}): wrap the AttnRes attention in gradient
        checkpointing.
\end{itemize}

\subsection{Integration with mHC and Mixture-of-Depths}

The residual strategy is \emph{orthogonal} to mHC and Mixture-of-Depths
(MoD): all combinations are supported. The interaction modes are:

\paragraph{mHC $\times$ AttnRes (\texttt{mhc\_stream\_mode}).}
When \texttt{use\_mhc=true} and the residual is \texttt{full\_attn} or
\texttt{block\_attn}, the depth-wise attention is applied either
\emph{per stream} (\texttt{"independent"}) or over the
\emph{flattened} $nC$-dim projection (\texttt{"joint"}). Independent
mode runs $n$ parallel attentions, one per stream --- matching the
paper's framing where each stream is a parallel "view" of the residual.
Joint mode treats the $n$ streams as a single fat residual, more
expressive but more parameter-hungry. The mHC coefficient projections
($\varphi_l$, $b_l$, $\alpha$ gates) are unaffected; AttnRes adds its
own parameters on top.

\paragraph{Mixture-of-Depths $\times$ AttnRes.}
MoD token routing operates \emph{within} a layer (selects which tokens
get the full layer update) and is orthogonal to the residual choice.
When \texttt{use\_mixture\_of\_depths=true} together with an AttnRes
residual, only the tokens selected by MoD contribute to the layer-output
buffer (Full AttnRes) or the intra-block partial sum (Block AttnRes);
the rest pass through unchanged and are scattered back by MoD's
existing logic. MoD is mutually exclusive with mHC (a
\texttt{ValueError} is raised in \texttt{HybridLayer.\-\_\-init\_\_} if
both are enabled); the corresponding combinations are
$(\text{no\_mHC}, \text{MoD}, \text{AttnRes})$,
$(\text{mHC}, \text{no\_MoD}, \text{AttnRes})$, or
$(\text{mHC}, \text{no\_MoD}, \text{no\_AttnRes})$.

\subsection{Reported Results}

The paper \citep{sun_attnres_2026} validates AttnRes through scaling
laws, ablations, and a 48B-parameter Kimi Linear extension:
\begin{itemize}[leftmargin=*]
    \item \textbf{Scaling laws}: both Full and Block AttnRes
        consistently outperform the PreNorm baseline across five model
        sizes. At 5.6 PFLOP/s-days, Block AttnRes reaches loss $1.692$
        versus the baseline's $1.714$ --- equivalent to a
        $1.25 \times$ compute advantage. The gap between Full and
        Block AttnRes narrows with scale (Full $1.737$ vs.\ Block
        $1.746$ on the 16-layer model).
    \item \textbf{Ablations}: input-dependent softmax attention
        ($1.737$) outperforms input-independent scalar mixing ($1.749$)
        and sigmoid attention ($1.741$). RMSNorm on keys is essential
        (Full $1.737$ vs.\ $1.743$ without; Block $1.746$ vs.\
        $1.750$). Multi-head depth aggregation
        ($\texttt{H} = 16$) hurts Block AttnRes ($1.752$ vs.\ $1.746$),
        indicating the optimal mixture is largely uniform across
        channels. Sliding-window attention over the last $W = 8$
        layers falls short of Block AttnRes ($1.764$ vs.\ $1.746$).
    \item \textbf{Kimi Linear extension}: a 48B-parameter
        (3B activated) Kimi Linear model with Block AttnRes
        ($N \approx 9$, $S = 6$) pre-trained on 1.4T tokens improves
        over the baseline across all 14 downstream benchmarks
        (MMLU-Pro, GPQA-Diamond, BBH, HellaSwag, TriviaQA, GSM8K,
        MGSM, Math, CMath, HumanEval, MBPP, CMMLU, C-Eval, ARC-C).
        The largest gains are on multi-step reasoning
        (GPQA-Diamond $+7.5$, Minerva Math $+3.6$) and code generation
        (HumanEval $+3.1$). Training-dynamics analysis shows that
        AttnRes mitigates PreNorm dilution: hidden-state magnitudes
        remain bounded across depth and gradient norms distribute more
        uniformly across layers.
    \item \textbf{Overhead}: training overhead is marginal when
        pipeline parallelism is not enabled; with pipeline parallelism
        Block AttnRes adds $<4\%$ end-to-end training overhead
        (cross-stage caching eliminates redundant transfers) and
        $<2\%$ inference latency overhead (two-phase computation with
        online softmax merging).
\end{itemize}

Example configurations: \texttt{configs/examples/full\_attn\_res\_adamw.yaml},
\texttt{configs/examples/block\_attn\_res\_adamw.yaml},
\texttt{configs/examples/mhc\_full\_attn\_res\_adamw.yaml}.
